% ---- pg-358 （书页 346） ----

\begin{mvframe}[显示板 8.1(续)]
{\ttfamily\small
\begin{center}
Principal Components Analysis
\end{center}
24 Observations\hfill\textbf{\rmfamily OUTPUT}\\
3 Variables

\vspace{12pt}
\begin{center}
Simple Statistics

\vspace{4pt}
\begin{tabular}{lccc}
 & X1 & X2 & X3\\
Mean & 4.725 443 647 & 4.477 573 765 & 3.703 185 794\\
StD & 0.105 223 590 & 0.080 104 466 & 0.082 296 771
\end{tabular}
\end{center}

\vspace{12pt}
\begin{center}
\fbox{Covariance Matrix}

\vspace{6pt}
\begin{tabular}{lccc}
 & X1 & X2 & X3\\
X1 & 0.011 072 004 0 & 0.008 019 141 9 & 0.008 159 648 0\\
X2 & 0.008 019 141 9 & 0.006 416 725 5 & 0.006 005 270 7\\
X3 & 0.008 159 648 0 & 0.006 005 270 7 & 0.006 772 758 5
\end{tabular}
\end{center}

\vspace{12pt}
\begin{center}
Total Variance = 0.024 261 488

\vspace{12pt}
\fbox{Eigenvalues of the Covariance Matrix}

\vspace{6pt}
\begin{tabular}{lcccc}
 & Eigenvalue & Difference & Proportion & Cumulative\\
PRIN1 & 0.023 303 & 0.022 705 & 0.960 508 & 0.960 51\\
PRIN2 & 0.000 598 & 0.000 238 & 0.024 661 & 0.985 17\\
PRIN3 & 0.000 360 & & 0.014 832 & 1.000 00
\end{tabular}

\vspace{12pt}
\fbox{Eigenvectors}

\vspace{6pt}
\begin{tabular}{lccc}
 & PRIN1 & PRIN2 & PRIN3\\
X1 & 0.683 102 & -.159 479 & -.712 697\\
X2 & 0.510 220 & -.594 012 & 0.621 953\\
X3 & 0.522 539 & 0.788 490 & 0.324 401
\end{tabular}
\end{center}
}
\end{mvframe}

\hyperref[fig:8.3]{图 8.3} 画出了崖底碎石图. 图上很突出的拐弯发生在 $i=2$ 处. 很清楚, 有一个支配地位的主成分.

解释 $96\%$ 总方差的第一主成分, 还有一个有趣的题材解释. 由于
